\clearpage
\section{Light Up (Akari)}
\subsection{Giới thiệu}
Light Up (tên gốc Akari) là trò chơi đặt bóng đèn trên lưới. Nhóm lấy luật từ
Nikoli và trang Puzzle Light Up~\cite{akari-rules,lightup}. Bàn gồm ô trắng và
ô tường màu đen; một số tường có ghi số từ 0 đến 4. Người chơi đặt bóng vào
các ô trắng sao cho:
\begin{itemize}
\item Mỗi bóng chiếu sáng theo hàng và cột cho tới khi gặp tường hoặc biên.
Ô có bóng cũng được tính là sáng.
\item Mọi ô trắng đều được chiếu sáng.
\item Hai bóng không được chiếu vào nhau. Hai bóng vẫn có thể cùng hàng hoặc
cùng cột nếu có tường chắn ở giữa.
\item Tường ghi số $k$ có đúng $k$ bóng ở bốn ô kề trên, dưới, trái, phải
(không tính đường chéo). Tường không ghi số chỉ có tác dụng chắn sáng.
\end{itemize}
Khác với Pipes, ở đây người chơi không xoay gì cả mà chọn ô để đặt bóng. Nhóm
chọn DFS làm giải thuật mù và backtracking có MRV và lan truyền ràng buộc
(gọi tắt là CSP) làm giải thuật heuristic.

\subsection{Phát biểu bài toán tìm kiếm}
Gọi $W$ là số ô trắng. Mỗi ô trắng $i$ có biến $x_i\in\{0,1\}$, với $x_i=1$ nghĩa
là ô đó có bóng. Chương trình lưu tập bóng, tập ô sáng và tập ô còn được phép
đặt bóng dưới dạng bitmask để kiểm tra nhanh.

\begin{table}[H]\centering\small
\begin{tabular}{@{}p{.20\textwidth}p{.74\textwidth}@{}}\toprule
Thành phần & Định nghĩa\\\midrule
Input & Vị trí ô trắng, tường và số trên tường\\
Trạng thái & Tập bóng đã đặt, tập ô đã sáng và tập ô còn được đặt bóng\\
Trạng thái đầu & Chưa có bóng nào\\
Hành động & DFS: đặt một bóng. CSP: đặt bóng tại một nguồn chiếu, hoặc chọn một tổ hợp bóng quanh tường\\
Chuyển trạng thái & Cập nhật bóng, ô sáng và ô còn đặt được; CSP suy luận thêm các bóng bắt buộc\\
Trạng thái đích & Mọi ô trắng sáng, không hai bóng chiếu nhau, mọi tường số đủ bóng\\
Chi phí & Số bóng của lời giải; cả hai giải thuật chỉ tìm lời giải đầu tiên, không tối ưu số bóng\\\bottomrule
\end{tabular}
\caption{Mô hình tìm kiếm của Light Up.}
\end{table}

Viết dưới dạng ràng buộc, gọi $V(i)$ là các ô trắng nhìn thấy từ $i$ (kể cả $i$),
$L$ là một đoạn ô trắng liền nhau trên một hàng hoặc cột, và $N(w)$ là các ô
trắng kề tường $w$ có số $k_w$:
\begin{align*}
\sum_{j\in V(i)}x_j &\geq 1 &&\text{(ô $i$ được chiếu sáng)},\\
\sum_{j\in L}x_j &\leq 1 &&\text{(không hai bóng chiếu nhau)},\\
\sum_{j\in N(w)}x_j &= k_w &&\text{(tường $w$ đủ bóng)}.
\end{align*}
Hàm \code{is\_goal\_state()} kiểm tra ba điều kiện này. Số tập bóng có thể có là
$2^W$; đây là kích thước không gian cấu hình, còn số trạng thái thực sự được
mở rộng nhỏ hơn rất nhiều nhờ cắt tỉa.

\begin{figure}[H]\centering
\begin{tikzpicture}[scale=.68]
\fill[orange!12] (-.5,.5) rectangle (6.5,-6.5);
\fill[ink!90] (2-.5,-0-.5) rectangle (2+.5,-0+.5);
\node[white,font=\footnotesize\bfseries] at (2,-0) {1};
\fill[ink!90] (3-.5,-0-.5) rectangle (3+.5,-0+.5);
\node[white,font=\footnotesize\bfseries] at (3,-0) {0};
\fill[ink!90] (5-.5,-1-.5) rectangle (5+.5,-1+.5);
\node[white,font=\footnotesize\bfseries] at (5,-1) {0};
\fill[ink!90] (0-.5,-3-.5) rectangle (0+.5,-3+.5);
\fill[ink!90] (4-.5,-3-.5) rectangle (4+.5,-3+.5);
\node[white,font=\footnotesize\bfseries] at (4,-3) {3};
\fill[ink!90] (1-.5,-4-.5) rectangle (1+.5,-4+.5);
\fill[ink!90] (5-.5,-5-.5) rectangle (5+.5,-5+.5);
\node[white,font=\footnotesize\bfseries] at (5,-5) {1};
\fill[ink!90] (1-.5,-6-.5) rectangle (1+.5,-6+.5);
\node[white,font=\footnotesize\bfseries] at (1,-6) {0};
\fill[orange!85] (0,-0) circle (.27);
\fill[orange!85] (6,-0) circle (.27);
\fill[orange!85] (2,-1) circle (.27);
\fill[orange!85] (4,-2) circle (.27);
\fill[orange!85] (1,-3) circle (.27);
\fill[orange!85] (5,-3) circle (.27);
\fill[orange!85] (4,-4) circle (.27);
\fill[orange!85] (0,-5) circle (.27);
\fill[orange!85] (5,-6) circle (.27);
\draw[ink!45] (-.5,.5) grid (6.5,-6.5);
\end{tikzpicture}

\caption{Bàn $7\times7$, mức easy, seed 70097 và lời giải do CSP tìm được.
Ô đen là tường, hình tròn là bóng đèn.}
\end{figure}

\subsection{Tạo dữ liệu đầu vào}
Bàn được sinh theo seed, nên cùng seed luôn cho cùng một bàn. Có ba mức
easy, normal và hard, khác nhau ở tỉ lệ tường $\rho$. Bàn $n\times n$ có
$M=\operatorname{round}(\rho n^2)$ ô tường:
\begin{center}\small
\begin{tabular}{@{}lrrrrr@{}}\toprule
Mức & $\rho$ & $7\times7$ & $10\times10$ & $14\times14$ & $25\times25$\\\midrule
easy   & 16\%   & 8 & 16 & 31 & 100\\
normal & 10.5\% & 5 & 10 & 21 & 66\\
hard   & 5.5\%  & 3 & 6 & 11 & 34\\
\bottomrule\end{tabular}
\end{center}
Các bước sinh bàn (\path{light_up_problem/generator/generate.py}):
\begin{enumerate}
\item Chọn ngẫu nhiên $M$ vị trí làm tường.
\item Lần lượt đặt bóng vào một ô trắng còn tối cho tới khi mọi ô đều sáng.
Tập bóng này là một lời giải mẫu.
\item Ghi số lên tường bằng số bóng kề trong lời giải mẫu, để khoảng 25\% số
tường không ghi số. Cuối cùng kiểm tra lại lời giải mẫu với luật.
\end{enumerate}
Vì bàn được dựng từ một lời giải nên chắc chắn có nghiệm. Solver không được
biết lời giải mẫu. Nhóm có tham khảo cách đánh giá độ khó theo kỹ thuật suy luận
của Dekker~\cite{akari-difficulty}, nhưng chưa áp dụng; nhãn easy/normal/hard
hiện chỉ phản ánh số tường, và chương trình chưa kiểm tra bàn có nghiệm duy nhất.

\subsection{Giải thuật tìm kiếm mù}
\subsubsection*{Depth-First Search (DFS)}
Ở mỗi bước, DFS lấy ô tối đầu tiên theo thứ tự hàng rồi cột và thử lần lượt
các vị trí có thể đặt bóng để chiếu ô đó. Một vị trí bị bỏ nếu bóng chiếu vào
bóng khác, làm một tường bị thừa bóng hay không còn đủ chỗ đặt bóng, hoặc làm
một ô tối không còn chỗ nào chiếu tới được. DFS đi sâu ngay vào nhánh hợp lệ
đầu tiên và quay lui khi bế tắc.

\noindent\begin{minipage}{\textwidth}
\begin{lstlisting}
DFS(state):
  if budget exhausted: return LIMIT
  if GOAL(state): return SOLVED(state)
  target = first dark cell in row-major order
  for bulb in legal illuminators(target):
    child = place(bulb)
    if FEASIBLE(child):
      result = DFS(child)
      if solved or limited: return result
    exclude bulb from later sibling choices
  return UNSOLVABLE
\end{lstlisting}
\end{minipage}

Dòng \code{exclude bulb} giúp các nhánh anh em không đặt lại một vị trí đã thử,
tránh việc cùng một tập bóng xuất hiện nhiều lần với thứ tự đặt khác nhau.
Trong code, đệ quy được thay bằng stack tường minh. Thứ tự chọn ô và thứ tự thử
là cố định, DFS không dùng điểm đánh giá nào.

Mỗi bước thêm một bóng nên độ sâu không quá $W$. Nếu mỗi ô tối có tối đa $b$
cách chiếu và lời giải sâu $d$, số trạng thái xấu nhất là $O(b^d)$. Stack chỉ
giữ nhánh hiện tại, nhưng chương trình còn lưu liên kết cha để phát lại các
bước, nên bộ nhớ tăng theo số trạng thái đã sinh.

\subsection{Giải thuật tìm kiếm heuristic}
\subsubsection*{Backtracking với MRV và lan truyền ràng buộc (CSP)}
Ý tưởng là suy ra những gì bắt buộc trước, chỉ đoán khi không suy ra được nữa,
và khi phải đoán thì đoán ở chỗ ít lựa chọn nhất. Cách làm này dựa trên phân
tích của Pulles~\cite{akari-analysis} và bộ giải Borroot/akari~\cite{akari-code}.

\paragraph{Trạng thái và tên các bước.} Dùng thống nhất
\code{state = (bulbs, lit, legal)}: tập bóng đã đặt, tập ô đã sáng và tập vị
trí còn được phép đặt bóng. Gọi $P$ là tập ô trắng, $dark=P\setminus lit$,
$V(i)$ là tập ô nhìn thấy từ $i$ (kể cả $i$), $N(w)$ là tập ô trắng kề tường
$w$, và $k_w$ là số ghi trên tường. Các tập này được lưu bằng bitmask trong code.
\code{PROPAGATE} là bước lan truyền ràng buộc;
\code{SELECT\_MRV\_CHOICES} chọn ràng buộc bằng MRV và trả các nhánh theo thứ
tự gain; \code{GOAL} kiểm tra đủ ba điều kiện thắng.

\paragraph{Lan truyền ràng buộc (CP): \code{PROPAGATE}.} Bước này lặp các luật
sau tới điểm cố định (không còn thay đổi), trước khi chọn nhánh:
\begin{itemize}
\item Tường đã đủ bóng thì cấm đặt bóng ở các ô kề còn lại. Tường còn thiếu $r$
bóng và chỉ còn đúng $r$ ô kề đặt được thì đặt bóng ở cả $r$ ô đó.
\item Ô tối chỉ còn một vị trí có thể chiếu tới thì đặt bóng ở vị trí đó.
Mỗi bóng mới làm cấm các ô khác trên đường sáng của nó.
\item Nếu đặt bóng ở một ô làm một tường số không thể đủ bóng, ô đó bị cấm.
\end{itemize}
Nếu gặp mâu thuẫn (ô tối không còn nguồn chiếu, tường không thể đủ bóng) thì
nhánh hiện tại bị bỏ.

\paragraph{Chọn ràng buộc bằng MRV: \code{SELECT\_MRV\_CHOICES}.}
Sau \code{PROPAGATE}, nếu chưa tới đích, MRV (Minimum Remaining Values) so
sánh số lựa chọn của hai loại ràng buộc trên cùng trạng thái:
\begin{itemize}
\item \textbf{Tường số $w$:} số bóng còn thiếu là
$r=k_w-|N(w)\cap bulbs|$; tập ô kề còn được phép đặt là
$U(w)=N(w)\cap legal$, có $u=|U(w)|$ phần tử. Với $r>0$, điểm là
\[
 m_{\mathrm{wall}}(w)=\binom{u}{r}.
\]
Mỗi lựa chọn đặt đúng $r$ bóng vào một tổ hợp của $U(w)$ và cấm các ô kề
không được chọn. Tường đã đủ bóng ($r=0$) không tham gia MRV; $r<0$ hoặc
$r>u$ là mâu thuẫn được phát hiện trong \code{PROPAGATE}.
\item \textbf{Ô tối $i$:} tập nguồn chiếu còn được phép là
$C(i)=V(i)\cap legal$. Điểm là
\[
 m_{\mathrm{dark}}(i)=|C(i)|.
\]
Mỗi lựa chọn đặt một bóng tại một nguồn chiếu trong $C(i)$; $C(i)=\varnothing$
là mâu thuẫn. Tập $legal$ đã loại các ô bị CP cấm, không chỉ xét ô có bóng hay chưa.
\end{itemize}
Chọn ràng buộc có điểm nhỏ nhất trong cả hai nhóm; bằng điểm thì ưu tiên
tường. Trong cùng loại, giữ phần tử gặp trước theo thứ tự duyệt của code.
Ví dụ, tường thiếu $r=2$ bóng và còn $u=3$ ô kề có điểm $3$; ô tối có hai
nguồn chiếu có điểm $2$, nên chọn ô tối. Nếu ô tối cũng có điểm $3$ thì chọn
tường. Điểm $\binom{u}{r}$ được tính \emph{trước} khi sinh và loại các tổ hợp
có bóng chiếu nhau, nên không nhất thiết bằng số nhánh thực tế còn lại.

\paragraph{Sắp thứ tự thử bằng gain.} Sau khi MRV chọn ràng buộc, gọi $A$ là
tập bóng mới của một nhánh: một bóng nếu chọn ô tối, hoặc $r$ bóng nếu chọn
tường. Điểm chiếu sáng thêm là
\[
 gain(A)=\left|\left(\bigcup_{b\in A}V(b)\right)\cap dark\right|
        =|lit_{\mathrm{child}}\setminus lit|.
\]
Lấy hợp trước khi đếm để một ô được nhiều bóng chiếu vẫn chỉ tính một lần.
Thử nhánh có gain lớn trước; bằng gain giữ thứ tự sinh ban đầu. Gain được tính
ngay sau khi đặt $A$, \emph{trước CP của trạng thái con}; không cộng gain vào
điểm MRV và không dùng nó để chọn giữa tường với ô tối.

Ví dụ thực từ bàn $7\times7$ easy, seed 70098: sau CP ở gốc, ô tối
$(6,7)$ có hai nguồn chiếu $(7,7)$ và $(5,7)$ (tọa độ hàng, cột tính từ 1).
Gain tương ứng là $7$ và $6$, nên thử $(7,7)$ trước
(Hình~\ref{fig:lightup-gain-70098}).

\begin{figure}[H]
\centering
\begingroup
\colorlet{blue}{pipeblue}
\colorlet{red}{reject}
\begin{subfigure}[t]{.46\textwidth}
\centering
\scalebox{1.55}{\begin{tikzpicture}[x=.35cm,y=-.35cm,font=\tiny]
% Snapshot after root CP: 7x7 easy, seed 70098; coordinates below are zero-based.
\foreach \c/\r in {0/0,2/0,3/0,4/0,5/0,6/0,2/1,3/1,4/1,6/1,0/2,1/2,3/2,4/2,5/2,0/3,1/3,2/3,3/3,4/3,5/3,6/3,0/4,2/4,5/4,0/5,2/5,0/6,2/6}{\fill[orange!12] (\c,\r) rectangle ++(1,1);}
\foreach \c/\r in {6/4,6/5,1/6,3/6,4/6,5/6,6/6}{\fill[blue!30] (\c,\r) rectangle ++(1,1);}
\foreach \c/\r/\n in {1/0/-1,0/1/2,1/1/1,5/1/3,2/2/2,6/2/-1,1/5/0,5/5/0}{
  \fill[ink] (\c,\r) rectangle ++(1,1);
  \ifnum\n>-1\node[white] at (\c+.5,\r+.5) {\n};\fi}
\foreach \i in {0,...,7}{
  \draw[muted!40,thin] (\i,0)--(\i,7);
  \draw[muted!40,thin] (0,\i)--(7,\i);}
\foreach \c/\r in {0/0,5/0,2/1,6/1,0/2,5/2,2/3}{\fill[orange!85!black] (\c+.5,\r+.5) circle (.16);}
\fill[blue] (6.5,6.5) circle (.22);
\draw[red!80!black,very thick] (6.08,5.08) rectangle (6.92,5.92);
\foreach \i in {0,...,6}{
  \pgfmathtruncatemacro{\labelnum}{\i+1}
  \node[muted] at (\i+.5,-.35) {\labelnum};
  \node[muted] at (-.35,\i+.5) {\labelnum};}
\end{tikzpicture}
}
\caption{Đặt bóng tại $(7,7)$: gain $=7$.}
\end{subfigure}\hfill
\begin{subfigure}[t]{.46\textwidth}
\centering
\scalebox{1.55}{\begin{tikzpicture}[x=.35cm,y=-.35cm,font=\tiny]
% Snapshot after root CP: 7x7 easy, seed 70098; coordinates below are zero-based.
\foreach \c/\r in {0/0,2/0,3/0,4/0,5/0,6/0,2/1,3/1,4/1,6/1,0/2,1/2,3/2,4/2,5/2,0/3,1/3,2/3,3/3,4/3,5/3,6/3,0/4,2/4,5/4,0/5,2/5,0/6,2/6}{\fill[orange!12] (\c,\r) rectangle ++(1,1);}
\foreach \c/\r in {1/4,3/4,4/4,6/4,6/5,6/6}{\fill[blue!30] (\c,\r) rectangle ++(1,1);}
\foreach \c/\r/\n in {1/0/-1,0/1/2,1/1/1,5/1/3,2/2/2,6/2/-1,1/5/0,5/5/0}{
  \fill[ink] (\c,\r) rectangle ++(1,1);
  \ifnum\n>-1\node[white] at (\c+.5,\r+.5) {\n};\fi}
\foreach \i in {0,...,7}{
  \draw[muted!40,thin] (\i,0)--(\i,7);
  \draw[muted!40,thin] (0,\i)--(7,\i);}
\foreach \c/\r in {0/0,5/0,2/1,6/1,0/2,5/2,2/3}{\fill[orange!85!black] (\c+.5,\r+.5) circle (.16);}
\fill[blue] (6.5,4.5) circle (.22);
\draw[red!80!black,very thick] (6.08,5.08) rectangle (6.92,5.92);
\foreach \i in {0,...,6}{
  \pgfmathtruncatemacro{\labelnum}{\i+1}
  \node[muted] at (\i+.5,-.35) {\labelnum};
  \node[muted] at (-.35,\i+.5) {\labelnum};}
\end{tikzpicture}
}
\caption{Đặt bóng tại $(5,7)$: gain $=6$.}
\end{subfigure}
\endgroup
\caption{Ví dụ thực từ bàn $7\times7$ easy, seed 70098, sau CP ở gốc và trước
CP của con. Viền đỏ đánh dấu ô tối $(6,7)$; chấm cam là bóng đã đặt, chấm xanh
là bóng thử. Nền cam nhạt là ô đã sáng, nền xanh là ô mới sáng dùng để tính gain.
Tọa độ hàng và cột tính từ 1.}
\label{fig:lightup-gain-70098}
\end{figure}

Khi chọn nguồn chiếu, các nhánh sau cấm các nguồn đã thử ở nhánh trước:
nhánh thứ hai đặt tại $(5,7)$ đồng thời cấm $(7,7)$. Các nhánh vì thế không
lặp cùng tập bóng theo thứ tự đặt khác nhau. Với tường, tổ hợp không hợp lệ
bị loại trước khi xếp gain; các mâu thuẫn còn lại được CP của con kiểm tra.

\paragraph{Quay lui: \code{BACKTRACK\_MRV\_CP}.} Mã giả dùng cùng tên các
bước ở trên. \code{SELECT\_MRV\_CHOICES} trả các trạng thái con đã áp dụng
lựa chọn và điều kiện cấm của từng nhánh, xếp theo gain giảm dần.

\noindent\begin{minipage}{\textwidth}
\begin{lstlisting}
BACKTRACK_MRV_CP(state):
  if budget exhausted: return LIMIT
  state, ok = PROPAGATE(state)
  if not ok: return UNSOLVABLE
  if GOAL(state): return SOLVED(state)
  checkpoint = state
  choices = SELECT_MRV_CHOICES(checkpoint)
  for child in choices:  # descending gain, before child CP
    result = BACKTRACK_MRV_CP(child)
    if result is SOLVED or LIMIT: return result
    state = checkpoint
  return UNSOLVABLE
\end{lstlisting}
\end{minipage}

Checkpoint giữ đủ \code{(bulbs, lit, legal)} sau CP của nút cha. Khi quay lui,
khôi phục cả ba tập trước khi thử nhánh kế tiếp. Code dùng stack tường minh
thay cho đệ quy; \code{PROPAGATE} tương ứng \code{propagate}, còn
\code{SELECT\_MRV\_CHOICES} tương ứng \code{select\_mrv\_choices} trong
\code{algorithms/search.py}. Mã giả gom các trường hợp dừng do tài nguyên
thành \code{LIMIT}; code kiểm tra giới hạn trong quá trình duyệt, và trong
CP khi bật trace.

Trường hợp xấu nhất CSP vẫn tốn thời gian hàm mũ. CP giúp giảm mạnh số nhánh
nhưng làm mỗi nút tốn công hơn, vì phải lặp các luật suy luận. Một nút CSP có
thể đặt nhiều bóng cùng lúc, nên số nút của CSP và DFS không tương ứng với cùng
một lượng công việc. Cả hai giải thuật chỉ bỏ những nhánh chắc chắn không có
nghiệm, nên sẽ tìm được lời giải nếu đủ thời gian.

\subsection{Hiện thực và minh họa từng bước}
Package \code{light\_up\_problem} gồm \code{core/} (luật và cấu trúc bàn),
\code{generator/} (sinh bàn), \code{algorithms/search.py} (giải thuật),
\code{ui/app.py} (giao diện pygame-ce và pygame-gui) và \code{benchmark/} (đo).
Solver chạy trong một worker riêng để giao diện không bị đứng.

\begin{itemize}
\item \textbf{Chơi tay:} click trái để đặt hoặc bỏ bóng, click phải để đánh dấu
ô không đặt bóng, rồi bấm kiểm tra lời giải.
\item \textbf{Xem giải thuật chạy:} chọn DFS hoặc CSP rồi bấm Start Search.
Có Pause/Resume, hai mũi tên để đi tới hoặc lui từng bước, và ba tốc độ
2, 8, 30 bước mỗi giây.
\item DFS hiện ô đang xét, bóng vừa đặt và lúc quay lui. CSP hiện thêm các bóng
được suy ra, các ô bị cấm (dấu $\times$) và việc khôi phục khi quay lui.
\item Nút \textbf{Compare DFS and CSP} chạy cả hai giải thuật một lần trên bàn
đang mở và hiện thời gian, bộ nhớ, số nút.
\end{itemize}

\subsection{Kết quả thực nghiệm}
Nhóm đo trên 4 kích thước, 3 mức và 5 seed cho mỗi cấu hình, tức 60 bàn khác
nhau. Mỗi bàn được giải một lần bằng mỗi giải thuật, tổng cộng 120 lượt. Mỗi
lượt chạy trong một process mới, các lượt chạy lần lượt và thứ tự hai giải thuật
được đảo ở các seed liên tiếp. Seed dùng cho từng cấu hình:
\begin{center}\small
\begin{tabular}{@{}rlll@{}}\toprule
Cỡ & easy & normal & hard\\\midrule
7 & 70097--70101 & 70194--70198 & 70291--70295\\
10 & 100097--100101 & 100194--100198 & 100291--100295\\
14 & 140097--140101 & 140194--140198 & 140291--140295\\
25 & 250097--250101 & 250194--250198 & 250291--250295\\
\bottomrule\end{tabular}
\end{center}

Mỗi lượt giới hạn 30 giây, 500.000 nút mở rộng và 512 MiB bộ nhớ. Lượt vượt
giới hạn được ghi là TIMEOUT, NODE\_LIMIT hoặc MEMORY\_LIMIT, không bị tính là
vô nghiệm. Thời gian đo bằng \code{perf\_counter} quanh solver, không tính thời
gian sinh bàn hay khởi động process. Bộ nhớ là đỉnh cấp phát của Python đo bằng
\code{tracemalloc}, không phải RAM của cả process. Mọi lời giải đều được kiểm tra
lại bằng \code{is\_goal\_state()}. Số liệu thô nằm ở
\path{docs/benchmarks/2026-10-31-lightup-dfs-csp.csv}.

\begin{center}\small
\begin{tabular}{@{}lrrrr@{}}
\toprule
Giải thuật & SOLVED & TIMEOUT & NODE\_LIMIT & MEMORY\_LIMIT\\
\midrule
DFS & 52 & 8 & 0 & 0\\
CSP & 60 & 0 & 0 & 0\\
\bottomrule
\end{tabular}
\end{center}
% Raw CSV: docs/benchmarks/2026-10-31-lightup-dfs-csp.csv.
% Arithmetic means; seconds/memory/nodes include TIMEOUT, bulbs use SOLVED only.
\begin{table}[H]
\centering\scriptsize
\setlength{\tabcolsep}{4pt}
\begin{tabular}{@{}rllcrrrr@{}}
\toprule
Cỡ & Mức & Giải thuật & Giải được & Thời gian (s) & Bộ nhớ (KiB) & Số nút & Số bóng\\
\midrule
7 & easy & DFS & 5/5 & 0.0102 & 27.1 & 30.6 & 8.4\\
7 & easy & CSP & 5/5 & 0.0078 & 25.1 & 4.0 & 8.4\\
\addlinespace[2pt]
7 & normal & DFS & 5/5 & 0.0067 & 26.8 & 9.2 & 7.6\\
7 & normal & CSP & 5/5 & 0.0103 & 29.1 & 6.2 & 7.2\\
\addlinespace[2pt]
7 & hard & DFS & 5/5 & 0.0072 & 27.3 & 8.4 & 7.4\\
7 & hard & CSP & 5/5 & 0.0105 & 30.5 & 6.4 & 6.2\\
\addlinespace[2pt]
10 & easy & DFS & 5/5 & 0.0360 & 58.7 & 51.6 & 15.4\\
10 & easy & CSP & 5/5 & 0.0280 & 56.0 & 7.8 & 13.8\\
\addlinespace[2pt]
10 & normal & DFS & 5/5 & 0.0280 & 62.6 & 40.2 & 13.0\\
10 & normal & CSP & 5/5 & 0.0312 & 62.2 & 8.8 & 11.6\\
\addlinespace[2pt]
10 & hard & DFS & 5/5 & 0.0230 & 62.4 & 12.4 & 11.4\\
10 & hard & CSP & 5/5 & 0.0314 & 67.7 & 9.8 & 10.0\\
\addlinespace[2pt]
14 & easy & DFS & 5/5 & 0.9070 & 228.9 & 983.4 & 26.2\\
14 & easy & CSP & 5/5 & 0.0947 & 108.3 & 12.4 & 24.4\\
\addlinespace[2pt]
14 & normal & DFS & 5/5 & 1.3044 & 375.7 & 2\,020.2 & 21.6\\
14 & normal & CSP & 5/5 & 0.1033 & 123.7 & 16.0 & 20.0\\
\addlinespace[2pt]
14 & hard & DFS & 5/5 & 0.0858 & 126.2 & 24.0 & 18.6\\
14 & hard & CSP & 5/5 & 0.1177 & 134.3 & 14.6 & 15.2\\
\addlinespace[2pt]
25 & easy & DFS & 2/5 & 23.7947 & 2192.9 & 8\,678.4 & 73.0\\
25 & easy & CSP & 5/5 & 2.1056 & 579.9 & 96.2 & 67.8\\
\addlinespace[2pt]
25 & normal & DFS & 1/5 & 24.2972 & 3522.0 & 13\,406.6 & 58.0\\
25 & normal & CSP & 5/5 & 1.0444 & 693.0 & 33.8 & 52.6\\
\addlinespace[2pt]
25 & hard & DFS & 4/5 & 7.7500 & 1268.7 & 2\,407.8 & 41.0\\
25 & hard & CSP & 5/5 & 1.0463 & 872.3 & 33.6 & 35.2\\
\bottomrule
\end{tabular}
\caption{Light Up: trung bình cộng (mean) trên 5 seed của mỗi cấu hình.
Thời gian, đỉnh bộ nhớ Python và số nút tính cả lượt TIMEOUT: thời gian là
thời lượng chạy tới lúc dừng, không phải thời gian giải xong. Số bóng chỉ lấy
trung bình các lượt SOLVED, với số mẫu ở cột Giải được. Kết quả từng bàn ở phụ lục.}
\label{tab:lightup-dfs-csp-vi}
\end{table}

\subsection{Nhận xét}
\paragraph{Số bàn giải được và số nút.} CSP giải được 60/60 bàn, DFS giải
được 52/60. Tám lượt TIMEOUT đều của DFS, ở bàn $25\times25$; không lượt nào
chạm giới hạn nút hay bộ nhớ. Tại $25\times25$ normal, DFS mở trung bình
13\,406.6 nút so với 33.8 nút của CSP. Ở cỡ 25, CSP mở trung bình 33.6--96.2
nút tùy nhãn. CP có thể suy ra nhiều bóng trong một nút, nên các số này không
đại diện cùng lượng công việc giữa hai giải thuật.

\paragraph{Thời gian.} Trong 52 bàn cả hai cùng giải được, CSP nhanh hơn ở
15 bàn, DFS nhanh hơn ở 37 bàn. Đây là số đếm theo từng cặp, khác với trung
bình theo cấu hình trong Bảng~\ref{tab:lightup-dfs-csp-vi}. Với bàn nhỏ,
chi phí lặp các luật CP có thể khiến CSP chậm hơn DFS; tuy nhiên CSP có mean
thời gian thấp hơn DFS ở $7\times7$ easy và $10\times10$ easy.
Tại $14\times14$ easy, mean DFS/CSP là 0.9070/0.0947 giây, tỉ số khoảng 9.6;
ở normal là 1.3044/0.1033 giây, tỉ số khoảng 12.6. Một số lượt DFS dài làm
mean tăng rõ rệt, nên không suy rộng các tỉ số này cho mọi seed.
Ở $25\times25$, CSP có mean 1.0444--2.1056 giây; mean thời lượng lượt chạy DFS
là 7.7500--24.2972 giây, có tính các lượt bị dừng sau 30 giây. Vì có TIMEOUT,
tỉ số thời lượng trung bình ở các nhóm này không phải tỉ số thời gian giải
xong trên cùng tập bàn.

\paragraph{Bộ nhớ và số bóng.} Ở cỡ 7 và 10, mean đỉnh bộ nhớ Python của
hai giải thuật khá gần nhau. Tại $25\times25$ easy và normal, mean DFS/CSP
lần lượt là 2192.9/579.9 KiB và 3522.0/693.0 KiB, tức khoảng 3.8 và 5.1 lần;
ở hard tỉ số khoảng 1.5 lần. DFS phải giữ nhiều trạng thái đã sinh để phát
lại đường đi. Cột Số bóng là mean riêng trên các lượt SOLVED: khi số lượt giải
được khác nhau, hai mean có thể dựa trên tập seed khác nhau. Trong 52 cặp
cùng giải được, CSP dùng ít bóng hơn ở 39 bàn; cả hai không bảo đảm tối ưu số bóng.

\paragraph{Hạn chế.} Mỗi cấu hình mới có 5 seed; mean nhạy với các lượt chạy
dài và các lượt bị giới hạn. Nhãn độ khó chỉ dựa trên số tường; ở cỡ 14 và 25,
mean thời lượng DFS ở hard thấp hơn easy và normal. Chương trình chưa kiểm
tra tính duy nhất của lời giải và chưa đo riêng tác dụng của MRV, gain và CP.
Các lượt đo cũ ngày 25/10 dùng generator trước đó được lưu riêng,
không trộn vào bảng thực nghiệm mới.
